\chapter{Results}
% - results and discussion
% - compare models and hyperparameters and feature sets and windows and attacks (easy vs stealthy performance)
% - And how it compares to IT


% plot correlation between parameters and certain feature sets and window lengthS?

% if an attacker could spend some time in the network they could observe the packets and try to match the metatada of his attack as much as possible

% possible big difference between test and val because of the misconfigurations in validation data.

\section{Evaluation Overview}
% Briefly explain models, splits, metrics and how to read the results.
\textcolor{blue}{is this too much repetition?}


This chapter presents the experimental results of the evaluated anomaly detection
models on the constructed reconnaissance detection datasets. The goal of the
evaluation is to determine to what extent unsupervised models trained only on
benign conntrack-derived traffic can detect injected reconnaissance activity in
real OT background traffic. As described in Section~\ref{sec:experimental_configuration}, the configurations vary in feature set, aggregation window length, anomaly threshold, and model-specific parameters.
For k-NN, the search space is more limited than for K-Means and the autoencoder
because of its higher computational cost. The k-NN results are therefore
interpreted as a targeted evaluation of validation-selected feature sets rather
than as a full sweep over all candidate feature sets.


The experiments are evaluated in two stages. First, the validation split is used
to compare configurations and identify promising model, feature-set and parameter
combinations. These validation results are used for model selection and for
understanding the effect of design choices such as window size and feature
selection. The selected configurations are then evaluated on the test split, which
is kept separate from the selection process and is used to assess how well the
chosen configurations generalize to another period of network traffic.

The results are reported using both window-level and range-level metrics.
Window-level metrics treat each source-host window as an individual prediction
and therefore measure how accurately the model classifies separate time windows.
Range-level metrics instead group consecutive anomalous windows into alarm
ranges and compare these ranges with the true attack intervals. This provides a
more operational view of the results, since reconnaissance activity usually
unfolds over a period of time rather than in a single isolated window.

Both views are used throughout the discussion. Window-level F1 indicates how
precise the model is at the individual-window level, while range-level F1
indicates whether complete reconnaissance scenarios are detected as temporal
events and how many separate false alarm ranges are produced. A model may
therefore have a moderate window-level score while still detecting the relevant
attack interval, or a strong window-level score while producing fragmented
alarms.

Unless stated otherwise, the tables in this chapter report configurations selected
according to validation performance, followed by their corresponding test-set
performance. Differences between validation and test results are interpreted as an
indication of robustness and sensitivity to the underlying benign traffic period.



\section{Validation-Based Feature Set Selection}
\label{sec:validation_feature_selection}
% Show top validation feature sets, explain reduced k-NN search.

Before comparing the final model configurations on the test split, the validation results were used to identify the most promising feature sets. This step was especially important because the full feature-set search could not be repeated exhaustively for k-NN due to its higher computational cost, as discussed in Section~\ref{sec:knn_resource_constraints}. The selection was therefore based on the full K-Means and autoencoder validation sweeps, which covered the complete planned feature-set space. 

\definecolor{selectedrow}{RGB}{230,242,255}

\
\begin{table}[htbp]
\centering
\small
\renewcommand{\arraystretch}{1.25}
\setlength{\tabcolsep}{3.5pt}
\makebox[\textwidth][c]{%
\begin{tabularx}{1.05\textwidth}{>{\raggedright\arraybackslash}X l l l l}
\toprule
Feature set & \#configs & Mean RF1 & Mean F1 & Score \\
\midrule

\rowcolor{selectedrow}
\feat{golden\_minimal} 
& 896 & 0.224 & 0.261 & 0.242 \\

\rowcolor{selectedrow}
\feat{omni\_core\_clean} 
& 896 & 0.214 & 0.250 & 0.232 \\

\feat{adam\_entropy} 
& 336 & 0.183 & 0.225 & 0.204 \\

\rowcolor{selectedrow}
\feat{low\_slow\_clean} 
& 896 & 0.127 & 0.157 & 0.142 \\

\feat{T2\_stealth\_scan} 
& 896 & 0.099 & 0.116 & 0.107 \\

\feat{T3\_T4\_T5\_loud\_scan} 
& 896 & 0.116 & 0.085 & 0.101 \\

\feat{recon\_core\_clean\_temporal} 
& 896 & 0.075 & 0.090 & 0.083 \\

\feat{ablation\_entropy\_only} 
& 896 & 0.061 & 0.050 & 0.055 \\

\feat{recon\_core\_clean} 
& 896 & 0.050 & 0.059 & 0.054 \\

\feat{horizontal\_scan\_clean} 
& 336 & 0.072 & 0.035 & 0.054 \\

\feat{ablation\_asymmetry\_only} 
& 336 & 0.000 & 0.000 & 0.000 \\

\bottomrule

\end{tabularx}
}
\caption{Validation-based feature set ranking. RF1 denotes range-level F1, and \#configs is the number of validation configurations included for each feature set. The selection score is computed as the average of mean validation RF1 and mean validation window-level F1. Highlighted rows indicate the feature sets selected for the final cross-model comparison and targeted k-NN evaluation.}
\label{tab:validation_feature_set_ranking}
\end{table}



The goal of this selection was not only to find the highest individual validation score, but to identify feature sets that performed well across multiple reconnaissance scenarios. Feature sets that were designed for one specific attack type were therefore interpreted with caution, even when they achieved strong results on that specific attack. Instead, preference was given to feature sets with broader coverage across the evaluated attacks. To summarize the validation results across feature sets, a compact selection score is used. This score gives equal weight to the mean validation range F1 and the mean validation window-level F1. Table~\ref{tab:validation_feature_set_ranking} shows the resulting ranking. 


The number of configurations is not identical for all feature sets. The general-purpose feature sets were included in the full K-Means and autoencoder validation sweeps. In contrast, \feat{csv\_adam\_entropy} was evaluated as an ADAM-inspired k-NN baseline and therefore has fewer configurations. It will be further discussed in Section~\ref{sec:adam_baseline}. Similarly, \feat{csv\_horizontal\_scan\_clean} and \feat{csv\_ablation\_asymmetry\_only} were dropped earlier based on their weaker validation behavior, so their scores are based on a smaller set of configurations. 

The two strongest feature sets are \feat{csv\_golden\_minimal} and \feat{csv\_omni\_core\_clean}, which both achieve the highest combined validation scores. \feat{csv\_low\_slow\_clean} is also retained because it is the next strongest general-purpose feature set and is relevant for slower reconnaissance behaviour. Although \feat{csv\_T2\_stealth\_scan} achieves competitive scores in some cases, it is more closely tied to a specific attack type and is therefore not selected.

The three highest-scoring general feature sets thus are \feat{csv\_golden\_minimal}, \texttt{csv\_omni\_core\_clean}, and \texttt{csv\_low\_slow\_clean}. These are therefore retained for the final cross-model comparison and for the reduced k-NN evaluation.



% \subsection{ADAM}
% In addition to the feature sets constructed specifically for reconnaissance
% detection, an ADAM-inspired feature set was also evaluated. As introduced in Section~\ref{sec:literature_models}, this feature set was
% included because ADAM uses entropy-based traffic characteristics for anomaly
% detection in cyber-physical systems, and therefore provides a useful comparison
% point for evaluating whether a more compact entropy-oriented representation is
% sufficient for the present conntrack-based reconnaissance detection task.

% The ADAM-inspired feature set was evaluated using the same validation procedure
% as the other candidate feature sets. However, its validation performance was not
% competitive and hence dropped at an early stage, as shown in Table~\ref{tab:adam_feature_set_results}


% FIXXXXXXXXX TABLEEEEEEEEEEEEEEE

% These results suggest that the ADAM-inspired representation does not capture
% enough of the reconnaissance-specific behaviour present in the constructed
% dataset. In particular, the strongest selected feature sets combine several types
% of indicators, including destination spread, short-flow behaviour, failed or
% asymmetric communication, and temporal aggregation. The ADAM-inspired feature set
% is more limited in this respect, which may explain its weaker performance on the
% evaluated reconnaissance scenarios.

\section{Overall Model Performance}
% Best final configurations, validation + test table.
\label{sec:overall_model_performance}
The next step is to compare the best-performing configurations across the evaluated model families. Since many model, feature-set, window-size and threshold combinations were evaluated, the focus lies on the strongest configurations and on the main patterns that emerge from the validation results. The full top-ranked validation configurations are provided in Appendix~\ref{tab:appendix_top_val_f1} and Appendix~\ref{tab:appendix_top_val_rf1}, where the models are ranked separately by window-level F1 and range-level F1. The appendix rankings show that the highest window-level F1 scores are dominated by CPS-GUARD configurations with 60-second windows. These configurations achieve very high window-level F1 and also strong range-level scores. In contrast, several of the highest range-F1 configurations use 5-minute windows and obtain perfect or near-perfect range precision with high range recall.


Both window-level F1 and range-level F1 are considered because they emphasize different aspects of the detection task. Window-level F1, rewards configurations that classify individual source-host windows accurately. Range-level F1 instead evaluates wheter complete reconnaissance intervals are detected as alamr ranges. A model can therefore obtain a high range-level score even when its window-level recall is lower, as long as it detects teh attack interval with few false alarm ranges. Similarly, a model can obtain a high window-level F1 while still producing fragmented alarms that reduce its range-level F1.

The two metrics also differ in the number of units over which they are computed. This becomes visible in the results. Window-level F1 is calculated over individual windows. Larger aggregation windows produce fewer evaluation points than shorter windows, which means that each false positive or false negative has a larger effect on the final score. Range-level F1 is calculated over detected and true alarm ranges rather than over individual windows. This usually results in even fewer evaluation units, so a single missed attack range or a single additional false alarm range can strongly affect the range-level precision, recall and F1-score. For this reason, the underlying precision and recall values are inspected together with F1 rather than using the F1-score alone.

The rankings also show why F1 should not be interpreted in isolation. Several configurations achieve similar F1 or range F1 values while making different trade-offs between precision and recall. In particular, a high range F1 can be obtained either by detecting almost all attack ranges with some additional false alarm ranges, or by producing very few false alarm ranges while missing one or more attack ranges. Although both types of error are relevant, missed attack ranges are more problematic in the context of early reconnaissance detection. A small number of additional false alarm ranges can still be inspected by an analyst, whereas a false negative means that an entire reconnaissance interval may remain unnoticed. Therefore, when configurations have comparable F1 scores, preference is given to models with no, or fewer, range false negatives. However, false positives remain important for operational usability, so the main ranking still uses F1 rather than a recall-weighted \(F_{\beta}\)-score.





\section{Influence of Window Size}
% Temporal granularity trade-offs.

\section{Influence of Feature Set}
% What feature families worked and why.

\section{Detection Performance per Reconnaissance Scenario}
% Loud scans, T2, T1/T0, low-and-slow.

\section{Validation--Test Generalization}
% Discuss discrepancies and robustness.

\section{Operational Discussion}
% SOC usefulness, false alarms, recommended configuration.
% Optional short IT comparison here if you have it.

\section{Limitations}
% Injected attacks, limited environments, validation dependence, reduced k-NN sweep.





% \section{Overall Detection Performance}
% Compare best k-NN, K-Means and CPS-GUARD-style autoencoder configurations.

% \section{Influence of Window Size}
% Discuss 1 min, 5 min, 1 h, 5 h trade-offs.

% \section{Influence of Feature Set}
% Discuss which feature families were useful and whether compact sets were enough.

% \section{Detection Performance per Reconnaissance Scenario}
% Discuss T5/T4/T3, T2, T1/T0, low-and-slow.

% \section{Comparison with IT Traffic}
% Discuss whether results transfer outside OT.

% \section{Operational Discussion}
% Discuss deployment meaning, false alarms, SOC usefulness, recommended setup.

% \section{Limitations}
% Discuss injected attacks, limited environments, noisy benign data, validation dependence.